\cotefile{82}
\scanpage{65}
à la section \deleted{\unsure{0|$\bullet$}} \inserted{nulle $s_0$} de \deleted{$\unsure{\mathbb P^1|\mathbb P}$} $X_0=\mathbb P^1$ (la section $\infty$, correspondant à la droite $\mathcal O_S e_2$, a comme stabilisateur
\[
\widetilde B'_0=\left\{\begin{pmatrix}\lambda&\unsure{0|b}\\ \mu&\lambda^{-1}\end{pmatrix}\ \middle|\ 
\begin{aligned}\lambda&\ \text{de }\unsure{\mathbb G_m|G_m}\\ \mu&\ \text{de }\mathbb G_a\end{aligned}\right\} \tag*{(11)}
\]

Comme l'action de $G_0$ sur $X_0$ \deleted{définit} section $s_0$ de $X_0$ définit
\[
\begin{aligned}g&\longmapsto g.s_0\\G_0&\longrightarrow X_0\end{aligned}\tag*{(12)}
\]
qui par passage au quotient définit un \uline{iso}
\[
G_0/B_0\longrightarrow X_0\ ,\tag*{(13)}
\]
\deleted{on trouve} compatible avec les op. de $G_0$, opérant par translation à g. sur \unsure{le} premier \deleted{\unsure{m}}\unsure{mb|b}, on trouve, en prenant l'image inverse sur $S$ et en tordant l'iso obtenu par le torseur $R$, l'iso.
\[
R\wedge X_{0S}\simeq R/B_{0S}\tag*{(14)}
\]
qui établit l'interprétation (7).

Le foncteur allant des \inserted{chacune des huit} catégories $C_{1S}\ldots$ vers $\operatorname{Tors}(\mathrm{GP}(1)_S)$ est le foncteur composé de
% archive sheet number: 64
